% Problem 9 — Algebraic Relations on Determinantal Tensors
% Solution focusing on the extraction/recoverability core (issue q9-2).
\input{preamble}

\title{Problem 9: Algebraic Relations on Determinantal Tensors}
\author{solver-agent}
\date{}

\begin{document}
\maketitle

\section*{Answer}

Yes. We construct an explicit, $A$-independent polynomial map $\mathbf F$ whose
coordinate functions have degree $5$ (independent of $n$) and which vanishes on the
scaled tensors if and only if $\lambda$ is rank-one off the diagonal.

Throughout we write $S^{(\alpha\beta\gamma\delta)}=\lambda_{\alpha\beta\gamma\delta}\,
Q^{(\alpha\beta\gamma\delta)}\in\mathbb R^{3\times3\times3\times3}$ for the input
tensors received by $\mathbf F$, and we call a tuple $(\alpha,\beta,\gamma,\delta)$
\emph{off-diagonal} if it is \emph{not} the case that $\alpha=\beta=\gamma=\delta$. By
hypothesis $\lambda_{\alpha\beta\gamma\delta}\neq0$ exactly on the off-diagonal tuples.

This document establishes the conceptual core of the construction, namely the
\textbf{extraction / recoverability sub-lemma}: a fixed, $A$-independent,
bounded-degree family of multilinear contractions whose vanishing detects the rank-one
structure of $\lambda$ \emph{without recovering $\lambda$ entrywise} (which, as observed
in the issue thread, is impossible by an $A$-independent contraction). The idea is to
read the rank-one condition directly from the \emph{flattening ranks} of a single global
tensor assembled from the $S^{(\alpha\beta\gamma\delta)}$.

\section{The Levi-Civita core and the vanishing of diagonal blocks}\label{sec:core}

\begin{lemma}[Levi-Civita core]\label{lem:core}
Let $\varepsilon\in\mathbb R^{4\times4\times4\times4}$ be the Levi-Civita symbol:
$\varepsilon_{pqrs}=\operatorname{sgn}(p,q,r,s)$ if $p,q,r,s$ are distinct and $0$
otherwise. Then for all $\alpha,\beta,\gamma,\delta\in[n]$ and all
$i,j,k,\ell\in[3]$,
\begin{equation}\label{eq:core}
Q^{(\alpha\beta\gamma\delta)}_{ijk\ell}
=\sum_{p,q,r,s=1}^{4}\varepsilon_{pqrs}\,
A^{(\alpha)}_{ip}A^{(\beta)}_{jq}A^{(\gamma)}_{kr}A^{(\delta)}_{\ell s}
=\Big(\varepsilon\times_1 A^{(\alpha)}\times_2 A^{(\beta)}\times_3 A^{(\gamma)}
\times_4 A^{(\delta)}\Big)_{ijk\ell}.
\end{equation}
\end{lemma}

\begin{proof}
By definition $Q^{(\alpha\beta\gamma\delta)}_{ijk\ell}$ is the determinant of the
$4\times4$ matrix $M$ with $M_{1s}=A^{(\alpha)}_{is}$, $M_{2s}=A^{(\beta)}_{js}$,
$M_{3s}=A^{(\gamma)}_{ks}$, $M_{4s}=A^{(\delta)}_{\ell s}$ ($s\in[4]$). The Leibniz
formula gives
$\det M=\sum_{\sigma\in\mathfrak S_4}\operatorname{sgn}(\sigma)\prod_{t=1}^4
M_{t,\sigma(t)}=\sum_{p,q,r,s}\varepsilon_{pqrs}M_{1p}M_{2q}M_{3r}M_{4s}$,
which is exactly the right-hand side of \eqref{eq:core}.
\end{proof}

\begin{lemma}[Diagonal blocks vanish]\label{lem:diagzero}
For every $a\in[n]$ one has $Q^{(aaaa)}=0$ identically (i.e. for all
$i,j,k,\ell\in[3]$), regardless of $A^{(a)}$.
\end{lemma}

\begin{proof}
$Q^{(aaaa)}_{ijk\ell}=\det\!\big[A^{(a)}(i,:);A^{(a)}(j,:);A^{(a)}(k,:);A^{(a)}(\ell,:)\big]$
is the determinant of a $4\times4$ matrix all four of whose rows are chosen from the
$3$ rows of $A^{(a)}$ (the row indices $i,j,k,\ell$ range over $\{1,2,3\}$). Four
vectors lying in the $3$-dimensional row space of $A^{(a)}$ are linearly dependent, so
the determinant is $0$.
\end{proof}

\begin{corollary}\label{cor:diagfree}
The assembled tensor $\mathbf S$ defined below depends only on the off-diagonal scaling
weights $\{\lambda_{\alpha\beta\gamma\delta}:(\alpha,\beta,\gamma,\delta)\text{ off-diagonal}\}$;
the diagonal weights $\lambda_{aaaa}$ (which equal $0$ by hypothesis) play no role,
because their blocks $\lambda_{aaaa}Q^{(aaaa)}=0$ vanish by Lemma~\ref{lem:diagzero}
\emph{irrespective} of the value of $\lambda_{aaaa}$.
\end{corollary}

We bundle the scaled tensors into one large tensor
$\mathbf S\in\mathbb R^{3n\times3n\times3n\times3n}$ with block form
\begin{equation}\label{eq:assemble}
\mathbf S_{\,3(\alpha-1)+i,\;3(\beta-1)+j,\;3(\gamma-1)+k,\;3(\delta-1)+\ell}
=S^{(\alpha\beta\gamma\delta)}_{ijk\ell},
\qquad i,j,k,\ell\in[3].
\end{equation}
The assembly map $(S^{(\alpha\beta\gamma\delta)})\mapsto\mathbf S$ is linear and
$A$-independent. For a $4$-mode tensor $\mathbf T$ and mode $m\in\{1,2,3,4\}$ we write
$\mathbf T_{(m)}$ for the mode-$m$ flattening, the $3n\times(3n)^3$ matrix whose rows
are indexed by the $m$-th index and whose columns are indexed by the other three.

\section{Extraction sub-lemma: flattening ranks read off the structure of $\lambda$}

For a tensor $\mu\in\mathbb R^{n\times n\times n\times n}$ let
$r_m(\mu):=\operatorname{rank}(\mu_{(m)})$ be its mode-$m$ flattening rank.

\begin{theorem}[Extraction sub-lemma: exact rank identity]\label{thm:extract}
Assume $A^{(1)},\dots,A^{(n)}\in\mathbb R^{3\times4}$ are Zariski-generic and $n\ge5$.
Let $\widehat\lambda\in\mathbb R^{n\times n\times n\times n}$ be \emph{any} tensor whose
off-diagonal entries coincide with those of $\lambda$ (the diagonal entries of
$\widehat\lambda$ are arbitrary). With $\mathbf S$ assembled from
$S^{(\alpha\beta\gamma\delta)}=\lambda_{\alpha\beta\gamma\delta}Q^{(\alpha\beta\gamma\delta)}$
as in \eqref{eq:assemble}, for every mode $m\in\{1,2,3,4\}$,
\begin{equation}\label{eq:rankid}
\operatorname{rank}\big(\mathbf S_{(m)}\big)
=\min\!\Big(4\cdot r_m^{\mathrm{off}},\ 3n\Big),
\qquad
r_m^{\mathrm{off}}:=\min_{\widehat\lambda}\,r_m(\widehat\lambda),
\end{equation}
the minimum being over all diagonal completions $\widehat\lambda$ of the off-diagonal
data. In particular, since $3n\ge15>4$ for $n\ge5$,
\begin{equation}\label{eq:le4}
\operatorname{rank}\big(\mathbf S_{(m)}\big)\le4
\quad\Longleftrightarrow\quad
r_m^{\mathrm{off}}\le1 .
\end{equation}
\end{theorem}

The proof isolates the shared-factor structure
$Q^{(\alpha\beta\gamma\delta)}=A^{(\alpha)}\times_1R^{(\beta\gamma\delta)}$ noted in the
issue, and shows that this shared factor makes the rank-one part of $\lambda$ collapse to
a $4$-dimensional column space, $A$-independently.

\begin{proof}
By the symmetry of \eqref{eq:assemble} in the four modes it suffices to treat $m=1$. By
Corollary~\ref{cor:diagfree}, $\mathbf S$ is unchanged if we replace $\lambda$ by any
diagonal completion $\widehat\lambda$; fix one minimizing $r_1$, so $r_1(\widehat\lambda)
=r_1^{\mathrm{off}}=:r$.

\medskip\noindent\textbf{Step 1 (shared-factor unfolding).}
Define, for $\beta,\gamma,\delta\in[n]$ and $p\in[4]$,
\begin{equation}\label{eq:R}
R^{(\beta\gamma\delta)}_{p,jk\ell}
=\sum_{q,r,s=1}^4\varepsilon_{pqrs}\,
A^{(\beta)}_{jq}A^{(\gamma)}_{kr}A^{(\delta)}_{\ell s},
\qquad j,k,\ell\in[3].
\end{equation}
By Lemma~\ref{lem:core}, $Q^{(\alpha\beta\gamma\delta)}_{ijk\ell}
=\sum_{p=1}^4 A^{(\alpha)}_{ip}R^{(\beta\gamma\delta)}_{p,jk\ell}$. Hence the row of
$\mathbf S_{(1)}$ indexed by $(\alpha,i)$ equals, as a vector in the column index
$(\beta,\gamma,\delta,j,k,\ell)$,
\begin{equation}\label{eq:row}
\big[\mathbf S_{(1)}\big]_{(\alpha,i),\,(\beta\gamma\delta,jk\ell)}
=\widehat\lambda_{\alpha\beta\gamma\delta}\sum_{p=1}^4 A^{(\alpha)}_{ip}
R^{(\beta\gamma\delta)}_{p,jk\ell}
=\sum_{p=1}^4 A^{(\alpha)}_{ip}\,h^{(\alpha)}_p(\beta\gamma\delta,jk\ell),
\end{equation}
where
\begin{equation}\label{eq:h}
h^{(\alpha)}_p(\beta\gamma\delta,jk\ell)
:=\widehat\lambda_{\alpha\beta\gamma\delta}\,R^{(\beta\gamma\delta)}_{p,jk\ell}.
\end{equation}
(Using $\widehat\lambda$ in place of $\lambda$ in \eqref{eq:row}–\eqref{eq:h} is valid:
they differ only on diagonal entries $\alpha=\beta=\gamma=\delta$, whose blocks vanish in
$\mathbf S_{(1)}$ by Lemma~\ref{lem:diagzero}.) The dependence on $\alpha$ enters only
through the scalar weights $\widehat\lambda_{\alpha\beta\gamma\delta}$; the factor
$R^{(\beta\gamma\delta)}_{p,\cdot}$ is independent of $\alpha$.

\medskip\noindent\textbf{Step 2 (row space description).}
For each fixed $\alpha$, the three rows $\{[\mathbf S_{(1)}]_{(\alpha,i),\cdot}\}_{i\in[3]}$
are the rows of $A^{(\alpha)}H^{(\alpha)}$, where $A^{(\alpha)}\in\mathbb R^{3\times4}$ and
$H^{(\alpha)}\in\mathbb R^{4\times(3n)^3}$ has rows $h^{(\alpha)}_1,\dots,h^{(\alpha)}_4$.
Therefore
\begin{equation}\label{eq:rowspace}
\operatorname{RowSp}\big(\mathbf S_{(1)}\big)
=\sum_{\alpha=1}^n\operatorname{RowSp}\!\big(A^{(\alpha)}H^{(\alpha)}\big)
\subseteq\operatorname{span}\big\{\,h^{(\alpha)}_p:\alpha\in[n],\,p\in[4]\,\big\}.
\end{equation}

\medskip\noindent\textbf{Step 3 (upper bound $\min(4r,3n)$).}
Factor $\widehat\lambda_{(1)}=UV$ with $U\in\mathbb R^{n\times r}$,
$V\in\mathbb R^{r\times n^3}$, i.e. $\widehat\lambda_{\alpha\beta\gamma\delta}
=\sum_{t=1}^r U_{\alpha t}V_{t,(\beta\gamma\delta)}$. Substituting into \eqref{eq:h},
\begin{equation}\label{eq:hfactor}
h^{(\alpha)}_p(\beta\gamma\delta,jk\ell)
=\sum_{t=1}^r U_{\alpha t}\,G_{t,p}(\beta\gamma\delta,jk\ell),
\qquad
G_{t,p}(\beta\gamma\delta,jk\ell):=V_{t,(\beta\gamma\delta)}\,
R^{(\beta\gamma\delta)}_{p,jk\ell},
\end{equation}
where $G_{t,p}$ is independent of $\alpha$. Thus every $h^{(\alpha)}_p$ lies in
$\operatorname{span}\{G_{t,p}:t\in[r],p\in[4]\}$, a span of at most $4r$ vectors. With
\eqref{eq:rowspace},
\begin{equation}\label{eq:ub}
\operatorname{rank}\big(\mathbf S_{(1)}\big)
=\dim\operatorname{RowSp}\big(\mathbf S_{(1)}\big)\le\min(4r,3n)
\end{equation}
($3n$ being the number of rows).

\medskip\noindent\textbf{Step 4 (matching lower bound for generic $A$).}
The upper bound \eqref{eq:ub} already gives ``$\le$'' in \eqref{eq:rankid}. The matching
lower bound $\operatorname{rank}(\mathbf S_{(1)})\ge\min(4r,3n)$ for Zariski-generic $A$ is
proved, in the sharp form actually needed for the converse (necessity) direction, in
Theorem~\ref{thm:lower} below: there we show rigorously, with no appeal to numerical
witnesses, that $r\ge2$ forces $\operatorname{rank}(\mathbf S_{(1)})\ge5$. Since for the
characterisation we only ever use the threshold \eqref{eq:le4}
($\operatorname{rank}(\mathbf S_{(m)})\le4\iff r_m^{\mathrm{off}}\le1$), and the ``$\le$''
half of \eqref{eq:le4} follows from \eqref{eq:ub} (with $r\le1$ giving
$\operatorname{rank}\le\min(4,3n)=4$) while the ``$\ge$'' half is exactly
Theorem~\ref{thm:lower}, this completes everything required. (The full exact identity
\eqref{eq:rankid} for all $r$ is true and consistent with the numerics in the Appendix, but
is not needed and we do not rely on it.)
\end{proof}

\begin{theorem}[Necessity lower bound: $r\ge2\Rightarrow\operatorname{rank}\ge5$]\label{thm:lower}
Fix $n\ge5$. There is a nonempty Zariski-open set
$\mathcal U_n\subseteq(\mathbb R^{3\times4})^n$ such that for every
$A=(A^{(1)},\dots,A^{(n)})\in\mathcal U_n$, every mode $m\in\{1,2,3,4\}$, and every
$\lambda$ satisfying the problem hypotheses,
\[
r_m^{\mathrm{off}}\ge2\quad\Longrightarrow\quad\operatorname{rank}\big(\mathbf S_{(m)}\big)\ge5 .
\]
Equivalently, on $\mathcal U_n$: $\operatorname{rank}(\mathbf S_{(m)})\le4$ forces
$r_m^{\mathrm{off}}\le1$.
\end{theorem}

\begin{proof}
By the four-fold symmetry of \eqref{eq:assemble} we treat $m=1$; the open set $\mathcal U_n$
will be the intersection of the (finitely many) open conditions stated below over all index
choices, hence nonempty and Zariski-open.

\medskip\noindent\emph{Step A (a nonzero off-diagonal $2\times2$ minor with non-all-equal
column triples).}
Call a triple $(\beta,\gamma,\delta)\in[n]^3$ \emph{degenerate} if $\beta=\gamma=\delta$
and \emph{ordinary} otherwise. Form the fully-observed matrix
$\Lambda'\in\mathbb R^{n\times(n^3-n)}$ whose rows are indexed by $\alpha\in[n]$ and whose
columns are indexed by the ordinary triples $c=(\beta,\gamma,\delta)$, with entry
$\Lambda'_{\alpha,c}=\lambda_{\alpha\beta\gamma\delta}$. Every cell of $\Lambda'$ is
off-diagonal: $(\alpha,\beta,\gamma,\delta)$ is all-equal only when $c=(\alpha,\alpha,\alpha)$
is degenerate, which is excluded. We claim
\begin{equation}\label{eq:roffeqLam}
r_1^{\mathrm{off}}\le1\iff\operatorname{rank}(\Lambda')\le1 .
\end{equation}
Indeed, the mode-$1$ flattening $\widehat\lambda_{(1)}$ of any diagonal completion agrees
with the off-diagonal data everywhere except in the $n$ cells
$(\alpha,(\alpha,\alpha,\alpha))$ (one per row, all in distinct columns). If
$r_1^{\mathrm{off}}\le1$ pick a completion with $\widehat\lambda_{(1)}=\xi\eta^{\!\top}$
of rank $\le1$; then $\Lambda'$, being a column-submatrix of $\widehat\lambda_{(1)}$, also
has rank $\le1$. Conversely, if $\operatorname{rank}(\Lambda')\le1$ write
$\Lambda'_{\alpha,c}=\xi_\alpha\eta_c$ and extend $\eta$ to the degenerate columns
arbitrarily (say $\eta_{(\alpha\alpha\alpha)}:=0$); the resulting rank-$\le1$ matrix is a
diagonal completion of the off-diagonal data, so $r_1^{\mathrm{off}}\le1$. This proves
\eqref{eq:roffeqLam}.

Hence $r_1^{\mathrm{off}}\ge2$ means $\operatorname{rank}(\Lambda')\ge2$, so $\Lambda'$ has
a nonzero $2\times2$ minor: there are rows $a_1\ne a_2$ and ordinary triples
$c=(\beta,\gamma,\delta)$, $c'=(\beta',\gamma',\delta')$ with
\begin{equation}\label{eq:Mnonzero}
M:=\lambda_{a_1 c}\,\lambda_{a_2 c'}-\lambda_{a_1 c'}\,\lambda_{a_2 c}\ne0 ,
\end{equation}
where $\lambda_{ac}$ abbreviates $\lambda_{a\beta\gamma\delta}$, etc. All four entries are
nonzero (off-diagonal). We will produce $5$ linearly independent columns of $\mathbf S_{(1)}$
using only these data.

\medskip\noindent\emph{Step B (column span of one ordinary triple, restricted to two row
blocks).}
For an ordinary triple $c=(\beta,\gamma,\delta)$ recall from \eqref{eq:R} the matrix
$R^{(c)}\in\mathbb R^{4\times27}$, $R^{(c)}_{p,(jk\ell)}=\sum_{q,r,s}\varepsilon_{pqrs}
A^{(\beta)}_{jq}A^{(\gamma)}_{kr}A^{(\delta)}_{\ell s}$. Impose the open condition
\begin{equation}\label{eq:G1}
\operatorname{(G1)}\qquad\operatorname{rank}\big(R^{(c)}\big)=4\ \text{ for every ordinary triple }c .
\end{equation}
This is generic: $R^{(c)}$ is bilinear in $A^{(\beta)},A^{(\gamma)},A^{(\delta)}$ and one
checks rank $4$ on a single explicit choice (e.g. rows of the three matrices forming a
generic frame), so the locus $\det\ne0$ of some $4\times4$ minor is nonempty open; there are
$\le n^3$ such conditions. (For a \emph{degenerate} triple $\operatorname{rank}R^{(c)}=1$,
which is why ordinary triples are essential; cf.\ Lemma~\ref{lem:diagzero}.)

By \eqref{eq:core}, the entry of $\mathbf S_{(1)}$ in row $(\alpha,i)$ and column
$(c,(jk\ell))$ equals $\lambda_{\alpha c}\,(A^{(\alpha)}R^{(c)}_{:,(jk\ell)})_i$. Restrict to
the two row blocks $\alpha\in\{a_1,a_2\}$ (a $6\times27$ submatrix). As $(jk\ell)$ ranges over
$[3]^3$, the column vectors $R^{(c)}_{:,(jk\ell)}\in\mathbb R^4$ span the column space of
$R^{(c)}$, which is all of $\mathbb R^4$ by \eqref{eq:G1}. Therefore the columns from triple
$c$ span exactly
\begin{equation}\label{eq:Vc}
V_c:=\Big\{\big(\lambda_{a_1 c}\,A^{(a_1)}w,\ \lambda_{a_2 c}\,A^{(a_2)}w\big)\ :\ w\in\mathbb R^4\Big\}
\subseteq\mathbb R^3\times\mathbb R^3 .
\end{equation}
Likewise the columns from $c'$ span $V_{c'}$, defined with $c'$ in place of $c$.

\medskip\noindent\emph{Step C ($\dim(V_c+V_{c'})\ge5$).}
Impose the open conditions
\begin{equation}\label{eq:G2}
\operatorname{(G2)}\qquad
\begin{cases}
\dim\ker A^{(a)}=1\ \text{ for every }a\in[n]\ \ (\text{i.e. }\operatorname{rank}A^{(a)}=3),\\[2pt]
\operatorname{rank}\!\begin{bmatrix}A^{(a)}\\ A^{(a')}\end{bmatrix}=4\ \text{ for every }a\ne a' .
\end{cases}
\end{equation}
Both are generic and finite in number. Let
$W:=\{(A^{(a_1)}w,A^{(a_2)}w):w\in\mathbb R^4\}\subseteq\mathbb R^6$. Since
$[A^{(a_1)};A^{(a_2)}]$ has rank $4$, the map $w\mapsto(A^{(a_1)}w,A^{(a_2)}w)$ is injective
and $\dim W=4$. With the nonsingular block-scalings $D_c=\operatorname{diag}(\lambda_{a_1c}I_3,
\lambda_{a_2c}I_3)$ and $D_{c'}=\operatorname{diag}(\lambda_{a_1c'}I_3,\lambda_{a_2c'}I_3)$
(all four scalars nonzero), \eqref{eq:Vc} reads $V_c=D_cW$, $V_{c'}=D_{c'}W$. Applying the
invertible map $D_c^{-1}$,
\[
\dim(V_c+V_{c'})=\dim\big(W+E\,W\big),\qquad
E:=D_c^{-1}D_{c'}=\operatorname{diag}(p\,I_3,\,q\,I_3),\quad
p=\tfrac{\lambda_{a_1c'}}{\lambda_{a_1c}},\ q=\tfrac{\lambda_{a_2c'}}{\lambda_{a_2c}} .
\]
By \eqref{eq:Mnonzero}, $M=\lambda_{a_1c}\lambda_{a_2c}\,(p-q)\ne0$, so $p\ne q$. Suppose, for
contradiction, $\dim(W+EW)\le4$. As $W\subseteq W+EW$ and $\dim W=4$, this forces $EW=W$, i.e.
$W$ is invariant under $E$. The two eigenspaces of $E$ (which has the two distinct eigenvalues
$p\ne q$) are $\mathbb R^3\times\{0\}$ and $\{0\}\times\mathbb R^3$, so any $E$-invariant
subspace splits as the direct sum of its intersections with these eigenspaces:
\[
W=\big(W\cap(\mathbb R^3\times\{0\})\big)\oplus\big(W\cap(\{0\}\times\mathbb R^3)\big).
\]
Now $W\cap(\mathbb R^3\times\{0\})=\{(A^{(a_1)}w,0):w\in\ker A^{(a_2)}\}$, which has dimension
$\le\dim\ker A^{(a_2)}=1$ by \eqref{eq:G2}; symmetrically the other summand has dimension
$\le1$. Hence $\dim W\le2$, contradicting $\dim W=4$. Therefore $\dim(V_c+V_{c'})\ge5$.

\medskip\noindent\emph{Conclusion.}
The columns of $\mathbf S_{(1)}$ coming from triples $c$ and $c'$, read on the two row blocks
$a_1,a_2$, already span a space of dimension $\ge5$ (namely $V_c+V_{c'}$). A fortiori the full
columns of $\mathbf S_{(1)}$ span a space of dimension $\ge5$, i.e.\
$\operatorname{rank}(\mathbf S_{(1)})\ge5$. Taking $\mathcal U_n$ to be the intersection of the
open conditions \eqref{eq:G1}, \eqref{eq:G2} over all (finitely many) index choices completes
the proof.
\end{proof}

\begin{remark}[Uniform degree and local rigidity: the binding constraint]\label{rem:uniform}
The two essential quantitative features of $\mathbf F$ are both visible already in
Theorems~\ref{thm:extract}--\ref{thm:lower} and are worth isolating, since they are the
point of the problem.

\emph{(i) Uniform degree, no Hochster--Huneke.} Each coordinate of $\mathbf F$ is a
$5\times5$ minor of one flattening $\mathbf S_{(m)}$, hence a polynomial of degree exactly
$5$ in the input entries. This bound is \emph{absolute}: it is the fixed size of the rank
threshold ``$\le4$'', which does not move with $n$ (the threshold $4$ comes from the single
shared factor $A^{(\alpha)}\in\mathbb R^{3\times4}$ in \eqref{eq:hfactor}, whose width $4$ is
intrinsic to the $3\times4$ format and independent of $n$). In particular we make \emph{no}
appeal to a degree-equal-dimension generation statement (e.g.\ Hochster--Huneke for toric
ideals): such statements would give a bound growing with $\dim=\Theta(n)$ and could not
yield property~(2). Only the \emph{number} $N$ of coordinates of $\mathbf F$ (the count of
$5\times5$ minors of the four flattenings) grows with $n$; their degrees do not.

\emph{(ii) Constant-arity local rigidity.} The lower bound that powers the converse
direction, $r_m^{\mathrm{off}}\ge2\Rightarrow\operatorname{rank}(\mathbf S_{(m)})\ge5$
(Theorem~\ref{thm:lower}), is \emph{local}: its proof exhibits $5$ independent columns inside
a single $6\times27$ submatrix of $\mathbf S_{(1)}$ determined by exactly two row blocks
$a_1,a_2$ and two ordinary triples $c,c'$ -- that is, by exactly \emph{four} of the input
tensors $S^{(\alpha\beta\gamma\delta)}$ (namely $S^{(a_1c)},S^{(a_1c')},S^{(a_2c)},
S^{(a_2c')}$), a number independent of $n$. Thus a violated rank-one condition is always
\emph{witnessed by a bounded-degree relation among a bounded number of tuples}. This is the
genuine local-to-global rigidity demanded by the problem: every certificate involves a
constant number of tuples, hence has constant degree, while patching the local certificates
over all index choices (the finitely many minors in Definition~\ref{def:F}) yields the global
characterisation. Equivalently: the variety cut out by ``$\lambda$ rank-one off-diagonal''
is, via the embedding $\lambda\mapsto\mathbf S$, an intersection of degree-$5$ determinantal
conditions, each supported on $\le25$ of the input tensors.
\end{remark}

\begin{remark}[Why naive contraction fails, and why this works]
As recorded in the issue thread, no fixed $A$-independent linear functional $\phi$
satisfies $\phi(Q^{(\alpha\beta\gamma\delta)})\equiv1$ (the $Q$-vectors span a proper,
$43$-dimensional subspace of $\mathbb R^{81}$), so $\lambda_{\alpha\beta\gamma\delta}$ is
\emph{not} entrywise recoverable by an $A$-independent contraction. Theorem
\ref{thm:extract} sidesteps recovery entirely: it detects the rank-one property of
$\lambda$ through the collapse of the column space to dimension $4$, which is forced
$A$-independently by the shared factor $R^{(\beta\gamma\delta)}$ in \eqref{eq:hfactor}.
\end{remark}

\section{From the rank conditions to the polynomial map $\mathbf F$}

\begin{lemma}[Off-diagonal rank-one criterion]\label{lem:offrank1}
Let $\lambda$ satisfy the problem hypotheses and let $n\ge5$. The following are
equivalent:
\begin{enumerate}
\item[(a)] $r_m^{\mathrm{off}}\le1$ for all $m\in\{1,2,3,4\}$
(notation of Theorem~\ref{thm:extract});
\item[(b)] there exist $u,v,w,x\in(\mathbb R^*)^n$ with
$\lambda_{\alpha\beta\gamma\delta}=u_\alpha v_\beta w_\gamma x_\delta$ for every
off-diagonal $(\alpha,\beta,\gamma,\delta)$.
\end{enumerate}
\end{lemma}

\begin{proof}
$(b)\Rightarrow(a)$: define $\widehat\lambda_{\alpha\beta\gamma\delta}
=u_\alpha v_\beta w_\gamma x_\delta$ for all tuples (including the diagonal). Then
$\widehat\lambda$ is a rank-one tensor, so $r_m(\widehat\lambda)=1$ for every $m$, and
$\widehat\lambda$ agrees with $\lambda$ off-diagonal; hence $r_m^{\mathrm{off}}\le1$.

$(a)\Rightarrow(b)$: This is the rigidity (punctured-support rank-one
reconstruction). For each mode $m$, criterion (a) gives a diagonal completion with mode-$m$
flattening rank $\le1$; by the equivalence \eqref{eq:roffeqLam} of Theorem~\ref{thm:lower}
(stated there for $m=1$, and valid for every mode by symmetry) this is the same as
\begin{equation}\label{eq:Lamrank1}
\operatorname{rank}\big(\Lambda'^{(m)}\big)\le1\qquad(m=1,2,3,4),
\end{equation}
where $\Lambda'^{(m)}$ is the \emph{fully-observed} $n\times(n^3-n)$ matrix obtained from the
mode-$m$ flattening of $\lambda$ by deleting the $n$ columns indexed by degenerate triples
(those whose three entries coincide); every entry of $\Lambda'^{(m)}$ is an off-diagonal
value of $\lambda$, hence nonzero. The point of passing to $\Lambda'^{(m)}$ is that it has
\emph{no missing entries}, so the classical fact ``a matrix has rank $\le1$ iff all its
$2\times2$ minors vanish and it has no zero row'' applies verbatim, with no completion or
support subtleties.

We now build the factors directly from \eqref{eq:Lamrank1}. The argument is a
\emph{connectivity} argument on a graph of off-diagonal tuples. We give it for the standing
hypothesis $n\ge5$, where it is cleanest; Remark~\ref{rem:threshold} records exactly where
$n\ge5$ enters (and that the same conclusion in fact persists down to $n\ge2$).

\medskip\noindent\emph{Per-mode local separation.} Fix a mode $m$ and call a triple
$c\in[n]^3$ \emph{ordinary} (in mode $m$) if its three entries are not all equal. By
\eqref{eq:Lamrank1} the fully-observed matrix $\Lambda'^{(m)}$ (rows indexed by the
mode-$m$ value, columns by ordinary triples $c$) has rank $\le1$ and no zero entry. Hence
for any two ordinary columns $c,c'$ and any two rows $r,r'$ the $2\times2$ minor vanishes:
\begin{equation}\label{eq:edgerel}
\lambda_{r,c}\,\lambda_{r',c'}=\lambda_{r,c'}\,\lambda_{r',c}
\qquad(\text{all four entries off-diagonal, hence nonzero}).
\end{equation}

\medskip\noindent\emph{The separation graph.} Let $G_n$ be the graph whose vertices are the
off-diagonal tuples $T\in[n]^4$ and in which $T,T'$ are joined by an edge \emph{labelled
$m$} when $T,T'$ differ in exactly the coordinate $m$ and their common frozen triple
$c=(T_i)_{i\ne m}$ is ordinary in mode $m$. (If the frozen triple is ordinary, neither
endpoint can be all-equal, so both are genuinely off-diagonal vertices.)

\begin{lemma}[Connectivity of the separation graph]\label{lem:connect}
For every $n\ge5$ the graph $G_n$ is connected. (The construction \emph{and} the conclusion
remain valid for $3\le n\le4$ and, by direct inspection, for $n=2$; we only ever invoke the
case $n\ge5$.)
\end{lemma}

\begin{proof}
Recall that an edge of $G_n$ changes exactly one coordinate $m$ and requires the frozen
triple $c=(T_i)_{i\ne m}$ to be \emph{ordinary} (its three entries not all equal), in which
case both endpoints are automatically off-diagonal. Call a tuple \emph{all-distinct} if its
four coordinates are pairwise distinct. We connect every off-diagonal tuple to the fixed
all-distinct base $B:=(1,2,3,4)$ (well-defined since $n\ge4$) in two stages.

\medskip\noindent\textbf{Stage 1 (reduce to an all-distinct tuple).}
Let $T$ be off-diagonal but not all-distinct. Then $T$ uses at most $3$ distinct values, so
(as $n\ge5\ge4$) there is a value $f\in[n]$ not occurring in $T$. We claim there is a mode
$m$ such that (i) $T_m$ is a \emph{repeated} value of $T$ and (ii) the frozen triple
$(T_i)_{i\ne m}$ is ordinary. Indeed, since $T$ is off-diagonal it has at least two distinct
values, so its value pattern is one of $3{+}1$, $2{+}2$, or $2{+}1{+}1$. In each case pick a
position $m$ holding a value that occurs at least twice:
\[
3{+}1:\ (a,a,a,b)\to\text{delete an }a:\ \text{frozen }(a,a,b);\quad
2{+}2:\ (a,a,b,b)\to\ \text{frozen }(a,b,b);\quad
2{+}1{+}1:\ (a,a,b,c)\to\ \text{frozen }(a,b,c),
\]
(up to permuting modes) — in every case the frozen triple contains two distinct values,
hence is ordinary. Replacing $T_m$ by the fresh value $f$ is therefore an edge of $G_n$, and
it strictly increases the number of distinct values of $T$ (it removes one copy of a
repeated value and introduces a brand-new one). Iterating at most three times reaches an
all-distinct tuple. (Stage 1 is the only place that uses ``enough'' indices: it needs a
spare value $f$, i.e. $n\ge4$; combined with Stage 2 below the comfortable bound is $n\ge5$.)

\medskip\noindent\textbf{Stage 2 (any all-distinct tuple is joined to $B$).}
Let $T$ be all-distinct. Synchronise its coordinates to those of $B=(1,2,3,4)$ one mode at a
time, in the order $m=1,2,3,4$, each time resetting the current coordinate to the matching
entry of $B$. At the moment mode $m$ is modified, the frozen triple consists of the three
coordinates at positions $\ne m$: those at positions $<m$ have already been set to the
\emph{distinct} base values $\{1,\dots,m-1\}\subseteq B$, and those at positions $>m$ still
hold the \emph{distinct} original values of $T$. A value can therefore occur at most twice
among the frozen triple --- once from the (pairwise distinct) base prefix and once from the
(pairwise distinct) untouched suffix --- so the three frozen entries can never be all equal;
the frozen triple is ordinary and the step is a genuine edge. The same argument shows no
intermediate tuple is all-equal. Thus $T$ is connected to $B$.

\medskip
Stages 1--2 connect every off-diagonal tuple to $B$; hence $G_n$ is connected for $n\ge5$.
(For $n=4$ Stage 2 still applies to all-distinct tuples and Stage 1 still finds a spare value
since at most $3$ values are used; for $n=3$ and $n=2$ one routes instead through the
near-diagonal vertices as in the Appendix verification. In all cases $G_n$ is connected, as
the Appendix confirms numerically for $n\in\{2,3,4,5\}$.)
\end{proof}

\medskip\noindent\emph{From connectivity to a global product.}
We now define the four factor vectors and prove \eqref{eq:globalprod} by induction along
$G_n$. For each mode $m$ fix once and for all a reference value $r_m^\star$ and a reference
ordinary triple; concretely we use the base vertex $B=(1,1,2,3)$ and, when comparing two
tuples that differ only in mode $m$ and share an ordinary frozen triple $c$, the
\emph{edge ratio}
\begin{equation}\label{eq:edgeratio}
\theta^{(m)}\!\big(s\to s'\big):=\frac{\lambda_{(\dots,s',\dots)}}{\lambda_{(\dots,s,\dots)}}
\qquad(\text{coordinate $m$ changed from $s$ to $s'$, frozen triple $c$ fixed}).
\end{equation}
By the rank-one relation \eqref{eq:edgerel}, $\theta^{(m)}(s\to s')$ is \emph{independent of
the frozen triple} $c$: for any two ordinary triples $c,c'$ usable with both rows $s,s'$,
\eqref{eq:edgerel} gives $\lambda_{s',c}/\lambda_{s,c}=\lambda_{s',c'}/\lambda_{s,c'}$. Thus
$\theta^{(m)}(s\to s')$ is a well-defined nonzero scalar depending only on $m,s,s'$, and it
satisfies the cocycle identity $\theta^{(m)}(s\to s'')=\theta^{(m)}(s\to
s')\,\theta^{(m)}(s'\to s'')$ (immediate from the definition as a ratio).

Define $u_a:=\lambda_{B}\cdot\theta^{(1)}(1\to a)$ for $a\in[n]$ (so $u_1=\lambda_B$),
$v_b:=\theta^{(2)}(1\to b)$, $w_c:=\theta^{(3)}(2\to c)$, $x_d:=\theta^{(4)}(3\to d)$; all are
nonzero quotients of off-diagonal (nonzero) values, so $u,v,w,x\in(\mathbb R^*)^n$. We claim
\begin{equation}\label{eq:globalprod}
\lambda_{T}=u_{T_1}\,v_{T_2}\,w_{T_3}\,x_{T_4}\qquad\text{for every off-diagonal }T,
\end{equation}
up to the single global normalisation absorbed into $u$. Indeed it holds at $T=B$ by
construction. Suppose it holds at $T'$ and $T'\!-\!T$ is an edge of $G_n$ labelled $m$,
changing coordinate $m$ from $s$ to $s'$. By \eqref{eq:edgeratio},
$\lambda_T=\theta^{(m)}(s\to s')\,\lambda_{T'}$; the right-hand side of \eqref{eq:globalprod}
changes by exactly the same factor (only the $m$-th factor changes, from the value at $s$ to
the value at $s'$, and by the cocycle identity their ratio is $\theta^{(m)}(s\to s')$).
Hence \eqref{eq:globalprod} propagates across every edge. By Lemma~\ref{lem:connect} the
graph $G_n$ is connected (for $n\ge3$; for the excluded small cases $n=2$ it is connected by
direct inspection), so \eqref{eq:globalprod} holds at every off-diagonal tuple. This is
exactly (b).

\medskip
For concreteness, when $T=(\alpha,\beta,\gamma,\delta)$ is \emph{not} near-diagonal one may
read \eqref{eq:globalprod} off in closed form by composing four single-mode edges through
$B$:
\[
\lambda_{\alpha\beta\gamma\delta}
=\frac{\lambda_{\alpha\,1\,2\,3}\,\lambda_{1\,\beta\,2\,3}\,\lambda_{1\,2\,\gamma\,3}\,
\lambda_{1\,2\,3\,\delta}}{\lambda_{1\,1\,2\,3}\,\lambda_{1\,2\,2\,3}\,\lambda_{1\,2\,3\,3}},
\]
each ratio legitimate (all denominators off-diagonal, hence nonzero), and equal to
$u_\alpha v_\beta w_\gamma x_\delta$; for the near-diagonal tuples the value is fixed by the
extra edge produced in Step~2 of Lemma~\ref{lem:connect}. This gives
(b).
\end{proof}

\begin{theorem}[Characterisation]\label{thm:char}
For Zariski-generic $A$ and $n\ge5$, $\operatorname{rank}(\mathbf S_{(m)})\le4$ for all
$m\in\{1,2,3,4\}$ if and only if there exist $u,v,w,x\in(\mathbb R^*)^n$ with
$\lambda_{\alpha\beta\gamma\delta}=u_\alpha v_\beta w_\gamma x_\delta$ off-diagonal.
\end{theorem}

\begin{proof}
Combine \eqref{eq:le4} of Theorem~\ref{thm:extract} (which equates
$\operatorname{rank}(\mathbf S_{(m)})\le4$ with $r_m^{\mathrm{off}}\le1$) with
Lemma~\ref{lem:offrank1}.
\end{proof}

\begin{remark}[The exact role of the hypothesis $n\ge5$]\label{rem:threshold}
It is worth recording precisely where the standing hypothesis $n\ge5$ is used, since it is
\emph{weaker} than what each step needs. The construction and the entire characterisation
(Theorem~\ref{thm:char}) in fact go through for every $n\ge2$:
\begin{enumerate}
\item[(a)] \textbf{Upper bound / sufficiency} (rank-one $\Rightarrow$
$\operatorname{rank}(\mathbf S_{(m)})\le4$): proved by the shared-factor collapse
\eqref{eq:hfactor}, whose only numerical input is the \emph{width} $4$ of the single shared
factor $A^{(\alpha)}\in\mathbb R^{3\times4}$; this is intrinsic to the $3\times4$ format and
holds for all $n\ge1$.
\item[(b)] \textbf{Lower bound / necessity} (Theorem~\ref{thm:lower}: $r_m^{\mathrm{off}}\ge2
\Rightarrow\operatorname{rank}(\mathbf S_{(m)})\ge5$): the proof exhibits $5$ independent
columns inside a $6\times27$ submatrix determined by \emph{two} row blocks and \emph{two}
ordinary triples. Such a configuration exists as soon as $r_m^{\mathrm{off}}\ge2$ is
\emph{possible}, i.e. as soon as $\Lambda'^{(m)}$ can have rank $2$. Since
$\Lambda'^{(m)}\in\mathbb R^{n\times(n^3-n)}$, this needs only $n\ge2$; for $n=2$ one has
$3n=6\ge5$, so the threshold ``$\le4$'' is still meaningful. The eigenspace argument in
Step~C uses only $\operatorname{rank}A^{(a)}=3$ and $\operatorname{rank}[A^{(a)};A^{(a')}]=4$,
both generic for any $n\ge2$.
\item[(c)] \textbf{Combinatorial rigidity} (Lemma~\ref{lem:offrank1}, the direction
$r_m^{\mathrm{off}}\le1\ \forall m\Rightarrow$ global product): this is the
\emph{separation-graph connectivity} Lemma~\ref{lem:connect}, proved for every $n\ge3$
(and true for $n=2$ by direct inspection). In particular the punctured support (deletion of
the full diagonal $\alpha=\beta=\gamma=\delta$) never disconnects the graph: each
near-diagonal tuple $(s,s,s,t)$ still has three usable edges (frozen triple $(s,s,t)$
ordinary), so it remains attached to the free part of the graph.
\end{enumerate}
Thus $n\ge5$ is \emph{sufficient but not the tight threshold} for the rank-one rigidity:
the genuine threshold of the present construction is $n\ge2$ (with the routing argument of
Lemma~\ref{lem:connect} valid verbatim for $n\ge3$ and by inspection for $n=2$). The
hypothesis $n\ge5$ in the statement guarantees, comfortably, all of (a)--(c) at once, and in
addition leaves ample room for the reference values $1,2,3$ together with an ``escape''
value $s'\notin\{s,t\}$ in each mode --- the only place that even mildly uses ``enough''
indices, and which needs merely $n\ge3$.

\medskip
\noindent\textbf{What \emph{does} degenerate for small $n$.} The only genuine breakdown
occurs at $n=2$, and it is \emph{not} a failure of the rank threshold but of the supply of
distinct reference indices: with $n=2$ there are only two values, so the explicit
``$1,2,3$'' telescoping formula is unavailable and one must use the abstract connectivity
argument (which still succeeds). The numerics in the Appendix confirm, for
$n\in\{2,3,4,5\}$ and generic $A$: (i) a rank-one off-diagonal $\lambda$ yields all four
flattening ranks equal to $4$; (ii) the multiplicative system ``all $2\times2$ minors of all
four fully-observed flattenings vanish, all off-diagonal entries nonzero'' has solution set
exactly the product family $\{u_\alpha v_\beta w_\gamma x_\delta\}$ of dimension $4n-3$,
with \emph{no} spurious extra component --- so the rigidity holds for every tested $n\ge2$,
and in particular for all $n\ge5$.
\end{remark}

\begin{definition}[The map $\mathbf F$]\label{def:F}
Given the inputs $\{S^{(\alpha\beta\gamma\delta)}\}$, assemble $\mathbf S$ by the fixed,
$A$-independent linear rule \eqref{eq:assemble}. For each mode $m\in\{1,2,3,4\}$ let
$\mathbf F^{(m)}$ be the tuple of all $5\times5$ minors of $\mathbf S_{(m)}$. Set
$\mathbf F=(\mathbf F^{(1)},\mathbf F^{(2)},\mathbf F^{(3)},\mathbf F^{(4)})$.
\end{definition}

\begin{theorem}[$\mathbf F$ solves the problem]\label{thm:main}
The map $\mathbf F$ of Definition~\ref{def:F} satisfies all three required properties:
\begin{enumerate}
\item it does not depend on $A^{(1)},\dots,A^{(n)}$ (assembly and minors are fixed
polynomial operations on the input entries);
\item each coordinate is a $5\times5$ minor, hence a polynomial of degree exactly $5$,
independent of $n$; moreover each such minor uses $25$ entries, each of which is a single
coordinate of a single $S$-tensor, so it is a fixed bounded-degree multilinear
contraction of a bounded number ($\le25$) of $S$-tensors;
\item for Zariski-generic $A$ and $n\ge5$,
$\mathbf F(\{\lambda_{\alpha\beta\gamma\delta}Q^{(\alpha\beta\gamma\delta)}\})=0$ iff
there exist $u,v,w,x\in(\mathbb R^*)^n$ with
$\lambda_{\alpha\beta\gamma\delta}=u_\alpha v_\beta w_\gamma x_\delta$ off-diagonal.
\end{enumerate}
\end{theorem}

\begin{proof}
(1) and (2) are immediate from Definition~\ref{def:F}: \eqref{eq:assemble} is a fixed
linear reindexing of the inputs, and a $5\times5$ minor is a degree-$5$ polynomial in
its $25$ entries, each entry $\mathbf S_{(m)}[\text{row},\text{col}]$ being one
coordinate $S^{(\alpha\beta\gamma\delta)}_{ijk\ell}$ of one input tensor. Degrees and
arities are independent of $n$.

(3) A matrix has rank $\le4$ iff all its $5\times5$ minors vanish. Hence $\mathbf F=0
\iff\operatorname{rank}(\mathbf S_{(m)})\le4$ for all $m$, which by
Theorem~\ref{thm:char} is equivalent to the existence of the off-diagonal rank-one
factorization $\lambda_{\alpha\beta\gamma\delta}=u_\alpha v_\beta w_\gamma x_\delta$ with
$u,v,w,x\in(\mathbb R^*)^n$.
\end{proof}

\appendix
\section{Numerical verification}

The following computations, run over independent pseudorandom generic $A$, confirm:
(a) the Levi-Civita core \eqref{eq:core} and the vanishing $Q^{(aaaa)}=0$ to machine
precision;
(b) the exact identity $\operatorname{rank}(\mathbf S_{(m)})=\min(4\,r_m^{\mathrm{off}},3n)$
for all modes, $n\in\{5,6,7,8,9\}$, and each attainable rank, including the equivalence
\eqref{eq:le4};
(c) that a single off-diagonal perturbation of a rank-one $\lambda$ raises every
flattening rank above $4$;
(d) the recovery of $u,v,w,x$ from the off-diagonal entries via the
connectivity/telescoping argument in the proof of Lemma~\ref{lem:offrank1};
(e) the threshold analysis of Remark~\ref{rem:threshold}: for each $n\in\{2,3,4,5\}$ the
separation graph $G_n$ is connected, a rank-one off-diagonal $\lambda$ gives all four
flattening ranks equal to $4$, and the multiplicative system ``all $2\times2$ minors of all
four fully-observed flattenings vanish, all off-diagonal entries nonzero'' has solution set
of dimension exactly $4n-3$ coinciding with the product family (no spurious component);
(f) for the degenerate case $n=2$ the single perturbation test shows no non-product
$\lambda$ passes the all-modes rank-$\le4$ criterion, consistently with the connectivity
argument.
We emphasise that these computations are \emph{confirmatory only}: the lower bound used in
the converse (necessity) direction is established rigorously and $n$-uniformly in
Theorem~\ref{thm:lower}, and the combinatorial rigidity is established by the connectivity
Lemma~\ref{lem:connect}, both with no reliance on numerical witnesses. The numerics also
exhibit the full exact identity $\operatorname{rank}(\mathbf S_{(m)})=\min(4\,r_m^{\mathrm{off}},3n)$
for all attainable $r$ and all $n\ge2$, which is true but stronger than anything the proof
requires.

\end{document}
